\documentclass[12pt]{article}
\usepackage{amsmath,amssymb,amsthm}
\usepackage{hyperref}
\usepackage[margin=1in]{geometry}

\newtheorem{theorem}{Theorem}
\newtheorem{lemma}[theorem]{Lemma}
\newtheorem{corollary}[theorem]{Corollary}
\newtheorem{proposition}[theorem]{Proposition}
\theoremstyle{definition}
\newtheorem{definition}[theorem]{Definition}

\title{Problem 322: Binomial Coefficients Divisible by 10}
\author{Project Euler Solution}
\date{}

\begin{document}
\maketitle

\section{Problem Statement}
Let $T(m,n)$ count the binomial coefficients $\binom{i}{n}$ divisible by 10 for $n\le i<m$.  Given $T(10^9,10^7-10)=989697000$, find $T(10^{18},10^{12}-10)$.

\section{Inclusion--Exclusion}

\begin{theorem}[Lucas]\label{thm:lucas}
For a prime $p$, $\binom{i}{n}\not\equiv 0\pmod{p}$ if and only if every base-$p$ digit of $n$ is at most the corresponding base-$p$ digit of $i$.
\end{theorem}

\begin{proof}
By Lucas' theorem, $\binom{i}{n}\equiv\prod_j\binom{i_j}{n_j}\pmod{p}$.  Each factor $\binom{i_j}{n_j}$ is nonzero in $\mathbb{F}_p$ precisely when $i_j\ge n_j$.
\end{proof}

\begin{definition}
Let $f_p(m,n)=\#\{i\in[n,m):p\nmid\binom{i}{n}\}$ and $f_{10}(m,n)=\#\{i\in[n,m):\gcd(\binom{i}{n},10)=1\}$.
\end{definition}

\begin{theorem}[Inclusion--exclusion]\label{thm:ie}
Since $10=2\times 5$ with $\gcd(2,5)=1$:
\[
T(m,n)=(m-n)-f_2(m,n)-f_5(m,n)+f_{10}(m,n).
\]
\end{theorem}

\begin{proof}
Standard inclusion--exclusion on the events ``$2\mid\binom{i}{n}$'' and ``$5\mid\binom{i}{n}$''.
\end{proof}

\section{Single-Prime Digit DP}

\begin{proposition}\label{prop:single}
For a prime $p$, $f_p(m,n)$ is computed by digit DP on the base-$p$ representation of $m$, enforcing $i_j\ge n_j$ at each position $j$, with a single boolean ``tight'' state.
\end{proposition}

\begin{proof}
Theorem~\ref{thm:lucas} requires digit-wise domination.  The standard digit-DP technique processes digits from most significant to least significant, tracking whether the prefix equals that of $m$ (tight) or is strictly smaller (free), ensuring exactly the integers in $[0,m)$ are enumerated.
\end{proof}

\section{Joint Digit DP for \texorpdfstring{$f_{10}$}{f10}}

\begin{theorem}\label{thm:joint}
When $m=10^K$, $f_{10}(m,n)$ is computed by a DP in base 10 with state $(c_2,c_5)$ representing carry values.  At position $k$ with digit $d$:
\begin{align}
\text{base-2 digit}_k &= (c_2+d)\bmod 2, \\
\text{base-5 digit}_k &= (c_5+d\cdot 2^k)\bmod 5, \\
c_2' &= \lfloor(c_2+d\cdot 5^k)/2\rfloor, \\
c_5' &= \lfloor(c_5+d\cdot 2^k)/5\rfloor.
\end{align}
After all $K$ positions, the final carries must dominate the remaining high-order digits of $n$.
\end{theorem}

\begin{proof}
Since $10^k=2^k\cdot 5^k$, each base-10 digit $d$ at position $k$ contributes $d\cdot 5^k$ to the base-2 expansion and $d\cdot 2^k$ to the base-5 expansion.  The carry variables track the accumulated higher-order contributions.  Since $5^k$ is always odd, $(c_2+d)\bmod 2$ correctly gives the base-2 digit at position $k$.
\end{proof}

\section{Result}

\[
T(10^{18},\,10^{12}-10)=\boxed{999998760323313995}
\]

\section{Complexity}
$O(K\cdot|S|\cdot 10)$ where $K=18$ and $|S|\le 3\times 10^5$ is the number of reachable carry states.

\section{Answer}

\[
\boxed{999998760323313995}
\]

\end{document}
